%% kapa: exam=1 round=27 year=2016 period=1 form=A subject=civil_law subject_ko=민법 seq=9 number=9
%% source: Q-Net 감정평가사 기출문제 (cst00309) artlSeq=5169402 file=extracted/1008/5169402/제27회 감정평가사 1차 문제/제27회 감정평가사 1교시문제(A형).pdf
%% conversion: kordoc 4.24.51 -> tools/gen1.py (auto); review=proofread-sonnet; flags=none
\begin{question}{9}
의사표시에 있어서 증명책임에 관한 설명으로 옳지 않은 것은? (다툼이 있으면 판례에 따름)
\begin{choices}
  \item 통정허위표시에서 제3자의 악의는 그 허위표시의 무효를 주장하는 자가 증명하여야 한다.
  \item 사기에 의한 의사표시에서 제3자의 악의는 취소를 주장하는 자가 증명하여야 한다.
  \item 진의 아닌 의사표시에서 상대방이 진의 아님을 알았거나 과실로 이를 알지 못하였다는 것은 의사표시의 무효를 주장하는 자가 증명하여야 한다.
  \item 상대방에게 도달하여야 효력이 있는 의사표시를 보통우편의 방법으로 하였다면, 송달의 효력을 주장하는 자가 그 도달을 증명하여야 한다.
  \item 착오로 인한 의사표시에서 착오가 법률행위 내용의 중요부분에 관한 것이라는 점과 중대한 과실이 없었다는 점은 표의자가 증명하여야 한다.
\end{choices}
\end{question}
\answer{5}
